\documentclass{easybook}
\usepackage{listings}

\begin{document}
\author{XiaoTang}
\date{2026.1.30}
该软件为德州仪器开发的用于DLP4500开发板的跨平台控制软件，编写了跨平台的编译方式。
用于控制DLP4500开发板的显示功能，并相当于开发参考。
\section{所需库}
项目UI需要QT6,DLP4500的SDK部分需求HIDAPI作为依赖。
\section{项目结构}
项目构建工具基于CMAKE与VCPKG，项目结构如下：
\begin{itemize}
    \item CMakeLists.txt:项目构建文件
    \item vcpkg.json:vcpkg依赖文件
    \item dlp:DLP4500 SDK文件夹
    \item qtui:UI相关文件夹
    \item cmake/toolchain:交叉编译工具链文件
    \item docker:LINUX交叉编译容器相关文件夹
    \item compile-guide:编译指南文件夹
\end{itemize}
\section{编译工具}
\subsection{编译链}
\subsubsection{Windows}
Windows平台可使用MINGW或MSVC编译器进行编译。
MSVC 工具链需要安装Visual Studio 生成工具，并在安装过程中选择“使用C++的桌面开发”并需勾选“用于 Windows 的 C++ CMake 工具”和“vcpkg包管理器”。
MINGW工具链需要安装MSYS2,运行“MSYS2.exe”进入MSYS的终端并调用以下指令安装相关工具和库：

\begin{lstlisting}
pacman -S mingw-w64-ucrt-x86_64-toolchain
\end{lstlisting}
即完成编译链安装。

\subsubsection{Linux}
Linux平台可使用GCC编译器进行编译。直接通过apt包管理器安装GCC即可：

\begin{lstlisting}
sudo apt-get update
sudo apt-get install build-essential
\end{lstlisting}

\subsection{CMAKE}
CMake是Modern C++工程管理的事实标准，它提供了一个平台无关的构建系统，可以生成适用于不同平台和编译器的构建文件。
其主要描述C++工程里各个文件之间的依赖关系，编译选项，链接选项等，并通过CMakeLists.txt文件进行配置。
依据配置自动生成各个编译器的构建文件，如Makefile，Ninja，Visual Studio项目文件等。
\subsubsection{Windows}
CMake的安装可以通过以下步骤完成：
\begin{itemize}
    \item MSVC工具链：安装Visual Studio 生成工具，勾选对应选项。
    \item MINGW工具链：安装\href{https://www.msys2.org/}{MSYS2},运行“MSYS2.exe”进入MSYS的终端并调用以下指令安装：

\begin{lstlisting}
pacman -S mingw-w64-ucrt-x86_64-cmake
\end{lstlisting}

\end{itemize}
\subsubsection{Linux}
Linux平台可以通过包管理器安装CMake，例如在Ubuntu上可以使用以下命令：

\begin{lstlisting}
sudo apt-get update
sudo apt-get install cmake
\end{lstlisting}
或Snap包管理器：

\begin{lstlisting}
sudo snap install cmake --classic
\end{lstlisting}
APT提供的版本较为落后，如需安装最新版本的CMake，使用Snap包管理器或参考CMake官方网站提供的
\href{https://apt.kitware.com/}{安装指南}

\subsection{VCPKG}
VCPKG是一个C++库管理工具，可以简化C++项目的依赖管理和构建过程。它提供了一个统一的接口来安装和管理C++库，使得开发者可以轻松地集成第三方库到他们的项目中。
VCPKG支持多平台和多编译器，可以自动处理库的下载、构建和安装过程，极大地简化了C++项目的依赖管理。
本项目使用VCPKG来管理依赖库HIDAPI和QT6。
\subsubsection{Windows}
在Windows平台上安装VCPKG，MSVC工具链安装Visual Studio 生成工具时勾选“vcpkg包管理器”选项即可完成安装。
MINGW工具链运行以下命令安装Git：

\begin{lstlisting}
pacman -S git
\end{lstlisting}
Git安装完成后打开命令行工具并执行以下命令：

\begin{lstlisting}
git clone https://github.com/Microsoft/vcpkg.git
cd vcpkg
.\bootstrap-vcpkg.bat
\end{lstlisting}
即可完成安装，安装完成后通过系统环境变量将vcpkg添加到PATH中并新建
“VCPKG\_ROOT”环境变量指向vcpkg的安装路径即可使用vcpkg命令，或者在使用时指定vcpkg的路径。

\subsubsection{Linux}
在Linux平台上安装VCPKG，首先需要安装Git和CMake，Git可以通过以下命令安装：
\begin{lstlisting}
sudo apt-get update
sudo apt-get install git
\end{lstlisting}
然后执行以下命令：
\begin{lstlisting}
git clone https://github.com/Microsoft/vcpkg.git
cd vcpkg
./bootstrap-vcpkg.sh
\end{lstlisting}
\section{本地编译}
在本地编译前参照前面的步骤需要先安装好相关工具和库，这里使用的IDE为vscode，visual studio 2022同理。
\subsection{前置插件}
在vscode中需要安装C/C++ Extension Pack插件。
\subsection{编译步骤}
在vscode中打开文件夹，选择工程目录，

\section{交叉编译}
\section{容器准备}
在vscode页面进入该目录下的docker文件夹进行容器构建，构建完成后从创建新的容器后进入。进入后，将vscode的工作区切换到工程下，安装好相关插件后vscode下方状态栏出现生成即可开始编译
\section{打包安装}
下方点击cpack后即可打包，打包的结果在./build文件夹下.sh文件,拷贝到开发板上后chmod后运行即可安装完成
\end{document}